% =========================================================
% frontmatter/objective.tex
% =========================================================

\chapter*{Project Objective}
\thispagestyle{fancy}
\addcontentsline{toc}{chapter}{Project Objective}
\begingroup
\setlength{\emergencystretch}{2em}

The objective of this project is to establish a traceable antenna-engineering workflow
for a center-fed resonant wire dipole operating near 1~GHz. The work connects a
closed-form half-wave baseline to a parameterized CST model, demonstrates resonance
control through physical-length tuning, and verifies the final radiation pattern against
canonical dipole theory.

The project has four principal technical objectives:

\begin{enumerate}
    \item Calculate the free-space wavelength and first-order half-wave dipole length at
    the 1~GHz design frequency.

    \item Construct the finite-radius, center-fed dipole model in CST Studio Suite and
    document the conductor geometry, center gap, discrete-port reference, boundary
    condition, frequency sweep, and far-field monitors.

    \item Tune the total dipole length through documented 150~mm, 125~mm, and
    135.6~mm cases, and quantify the final resonant frequency and reflection coefficient.

    \item Compare the final CST far-field behavior with an independent analytical
    calculation of the thin half-wave dipole directivity and half-power beamwidth.
\end{enumerate}

A secondary objective is to distinguish exact marker values, displayed figure values,
and plot-derived estimates, while documenting the numerical and experimental evidence
still required before a hardware-performance claim would be justified.

The final outcome is intended to provide a concise, university-style engineering report
that remains suitable for a public technical portfolio. All reported performance results
are analytical or simulated; no fabricated prototype or measured data are claimed.
\endgroup
